\subsection{Spectral radius}

Lemma 2 (Fekete's Lemma). --- Let \((a_{n})_{n\geqslant1}\) be a sequence of real numbers. Suppose that

\[
a _ {n + m} \leqslant a _ {n} + a _ {m}
\]

for all \(n \geqslant 1\) and every \(m \geqslant 1\) . Then the sequence \((a_n / n)_{n \geqslant 1}\) converges and satisfies

\[
\lim _ {n \to + \infty} \frac {a _ {n}}{n} = \inf _ {n \geqslant 1} \frac {a _ {n}}{n}.
\]

Let us set \(a_0 = 0\) ; the inequality \(a_{n+m} \leqslant a_n + a_m\) remains valid for every \(n \geqslant 0\) and every \(m \geqslant 0\) . Fix an integer \(m \geqslant 1\) . For every integer \(n \geqslant 1\) , let \(q(n)\) and \(r(n)\) be the integers such that \(n = q(n)m + r(n)\) and \(0 \leqslant r(n) < m\) (E, III, p.~39, th. 1). The hypothesis then implies

\[
\frac {a _ {n}}{n} \leqslant \frac {a _ {q (n) m}}{n} + \frac {a _ {r (n)}}{n} \leqslant \frac {q (n) a _ {m}}{n} + \frac {a _ {r (n)}}{n} \leqslant \frac {q (n)}{n} a _ {m} + \frac {m}{n}.
\]

Letting \(n\) tend to \(+\infty\) , one deduces that \(\limsup_{n}(a_{n}/n) \leqslant a_{m}/m\) since \(q(n)/n \to 1/m\) . Since this holds for every \(m \geqslant 1\) , we therefore have

\[
\limsup _ {n \to + \infty} \frac {a _ {n}}{n} \leqslant \inf _ {m \geqslant 1} \frac {a _ {m}}{m} \leqslant \liminf _ {n \to + \infty} \frac {a _ {n}}{n}.
\]

These inequalities prove the convergence of the sequence \((a_{n}/n)_{n\geqslant1}\) as well as the formula \(\lim a_{n}/n = \inf_{n\geqslant1} a_{n}/n\) .

PROPOSITION 1. --- Let A be a normed algebra. For every \(x \in \mathrm{A}\) the sequence \((\|x^n\|^{1/n})_{n \geqslant 1}\) is convergent and its limit \(\varrho(x)\) is equal to \(\inf_{n \geqslant 1} \|x^n\|^{1/n}\) . Moreover, for every norm \(x \mapsto \|x\|_1\) defining the topology of A, one also has

\[
\varrho (x) = \lim _ {n \to + \infty} \| x ^ {n} \| _ {1} ^ {1 / n} = \inf _ {n \geqslant 1} \| x ^ {n} \| _ {1} ^ {1 / n}.
\]

If \(x\) is nilpotent, one has \(\| x^n\| _1^{1 / n} = 0\) for every integer \(n\) sufficiently large and every norm \(x\mapsto \| x\| _1\) defining the topology of A.

Suppose now that \(x\) is not nilpotent, and let us set \(\alpha_{n} = \| x^{n}\|\) . We have \(\alpha_{n} > 0\) for every integer \(n \geqslant 1\) , and \(\alpha_{n + m} \leqslant \alpha_{n}\alpha_{m}\) for all \(n, m \in \mathbf{N}\) by (1). Lemma 2, applied to the sequence \(a_n = \log(\alpha_n)\) shows the existence of the limit \(\varrho(x)\) and the formula \(\varrho(x) = \inf_{n > 0} \alpha_n^{1/n}\) .

Let \(x \mapsto \|x\|_1\) be a norm defining the topology of A. There exist real numbers \(a > 0\) and \(b > 0\) such that

\[
a \| x \| \leqslant \| x \| _ {1} \leqslant b \| x \|
\]

for all \(x \in A\) (EVT, II, p.~7, cor. 2). Consequently,

\[
a ^ {1 / n} \| x ^ {n} \| ^ {1 / n} \leqslant \| x ^ {n} \| _ {1} ^ {1 / n} \leqslant b ^ {1 / n} \| x ^ {n} \| ^ {1 / n},
\]

for all \(n \geqslant 1\) ; whence, passing to the limit, or taking the infimum, the equality

\[
\varrho (x) = \lim _ {n \to + \infty} \| x ^ {n} \| _ {1} ^ {1 / n} = \inf _ {n > 0} \| x ^ {n} \| _ {1} ^ {1 / n}.
\]

DEFINITION 2. --- For every element x of a normed algebra A, the real number

\[
\varrho (x) = \lim _ {n \to \infty} \| x ^ {n} \| ^ {1 / n} = \inf _ {n > 0} \| x ^ {n} \| ^ {1 / n}
\]

is called the spectral radius of x.

For every element \(x\) of A, we have

\[
\begin{array}{r l} (3) & \varrho (x) \leqslant \| x \| \\ (4) & \varrho (x ^ {n}) = \varrho (x) ^ {n}, \text { for every integer } n \geqslant 1. \end{array}
\]

DEFINITION 3. --- An element x of A is quasi-nilpotent if \(\varrho(x)=0\) .

This amounts to saying that, for any \(\lambda \in \mathbf{K}\) , the numbers \(\| (\lambda x)^n\|\) are bounded for \(n\geqslant 1\) , or equivalently that, for any \(\lambda \in \mathbf{K}\) , the sequence \((\lambda x)^n\) tends to 0 as \(n\to +\infty\) .

Remark 1. --- Let A be a normed algebra. If an element \(x \in \mathbf{A}\) satisfies \(\varrho(x) = \|x\|\) , one has \(\|x^n\| = \|x\|^n\) for all \(n \in \mathbf{N}\) , by (3) and (4).

Conversely, suppose that \(\|x^{2}\|=\|x\|^{2}\) for all \(x\in A\) . Then for every integer \(n\geqslant0\) , we have the equality \(\|x^{2^{n}}\|=\|x\|^{2^{n}}\) , hence \(\|x\|=\|x^{2^{n}}\|^{2^{-n}}\) ; as n tends to \(+\infty\) , we obtain \(\|x\|=\varrho(x)\) .

Remark 2. --- The function \(x \mapsto \varrho(x)\) on A, being the lower envelope of continuous functions \(x \mapsto \|x^n\|^{1/n}\) for \(n \geqslant 1\) , is upper semicontinuous (TG, IV, p.~31, cor.), but in general it is not continuous. It can even happen (cf.~exerc. 12 of I, p.~157) that a sequence of nilpotent elements of A tends to an element which is not quasi-nilpotent.

